% Part: computability
% Chapter: recursive-functions
% Section: pr-functions

\documentclass[../../../include/open-logic-section]{subfiles}

\begin{document}

\olfileid{cmp}{rec}{prf}
\olsection{الدوال العودية البدائية}

لنسجل مرة أخرى كيف يمكن تعريف دوال جديدة من دوال موجودة باستعمال
العودية البدائية والتركيب.

\begin{defn}
  \ollabel{defn:primitive-recursion} افترض أن $f$ دالة ذات $k$ حجج
  ($k\ge 1$) وأن $g$ دالة ذات $k+2$ حجة. تكون الدالة المعرفة
  \emph{بالعودية البدائية من $f$ و$g$} هي الدالة~$h$ ذات $k+1$ حجة
  التي تعرفها المعادلتان
  \begin{align*}
    h(x_0,\dots,x_{k-1},0) & =  f(x_0,\dots,x_{k-1}) \\
    h(x_0,\dots,x_{k-1},y+1) & = g(x_0,\dots,x_{k-1}, y, h(x_0,\dots,x_{k-1}, y))
  \end{align*}
\end{defn}

\begin{defn}
  \ollabel{defn:composition} افترض أن $f$ دالة ذات $k$ حجج، وأن
  $g_0$، …،~$g_{k-1}$ دوال عددها $k$ وكلها ذات $n$ حجة. تكون
  الدالة المعرفة \emph{بتركيب $f$ مع $g_0$، …،~$g_{k-1}$} هي الدالة
  ~ $h$ ذات $n$ حجة التي تعرفها
  \[
  h(x_0, \dots, x_{n-1}) =
  f(g_0(x_0, \dots, x_{n-1}), \dots, g_{k-1}(x_0, \dots, x_{n-1})).
  \]
\end{defn}

سنعد، إلى جانب $\Succ$ ودوال الإسقاط
\[
\Proj{n}{i}(x_0,\dots,x_{n-1}) = x_i,
\]
لكل عدد طبيعي $n$ ولكل $i<n$، الدالة~$\Zero(x)=0$ من الدوال العودية
البدائية.

\begin{defn}
  مجموعة الدوال العودية البدائية هي مجموعة الدوال من $\Nat^n$ إلى
  $\Nat$، المعرفة استقرائيًا بالبنود الآتية:
  \begin{enumerate}
  \item $\Zero$ عودية بدائية.
  \item $\Succ$ عودية بدائية.
  \item كل دالة إسقاط $\Proj{n}{i}$ عودية بدائية.
  \item إذا كانت $f$ دالة عودية بدائية ذات $k$ حجج وكانت $g_0$،
    …،~$g_{k-1}$ دوال عودية بدائية ذات $n$ حجة، فإن تركيب $f$ مع
    $g_0$، …،~$g_{k-1}$ عودي بدائي.
  \item إذا كانت $f$ دالة عودية بدائية ذات $k$ حجج وكانت $g$ دالة
    عودية بدائية ذات $k+2$ حجة، فإن الدالة المعرفة بالعودية البدائية
    من $f$ و$g$ عودية بدائية.
\end{enumerate}
\end{defn}

\begin{explain}
وبعبارة أوجز، مجموعة الدوال العودية البدائية هي أصغر مجموعة تحتوي
$\Zero$ و$\Succ$ ودوال الإسقاط~$\Proj{n}{j}$، وتكون منغلقة تحت
التركيب والعودية البدائية.

وثمة طريقة أخرى لوصف مجموعة الدوال العودية البدائية، هي تعريفها بدلالة
«مراحل». لتدل $S_0$ على مجموعة دوال البدء: $\Zero$ و$\Succ$ ودوال
الإسقاط. وهذه هي الدوال العودية البدائية من المرحلة~$0$. ومتى عرفت
مرحلة $S_i$، فلتكن $S_{i+1}$ مجموعة جميع الدوال الناتجة من تطبيق حالة
واحدة من التركيب أو العودية البدائية على دوال موجودة من قبل في~$S_i$.
وعندئذ تكون
\[
S = \bigcup_{i \in \Nat} S_i
\]
مجموعة جميع الدوال العودية البدائية.
\end{explain}

لنتحقق من أن $\Add$ دالة عودية بدائية.

\begin{prop}
  دالة الجمع $\Add(x,y)=x+y$ عودية بدائية.
\end{prop}

\begin{proof}
لدينا من قبل تعريف بالعودية البدائية لـ$\Add$ بدلالة دالتين $f$ و~$g$
يطابق صورة \olref{defn:primitive-recursion}:
\begin{align*}
  \Add(x_0, 0) & = f(x_0) = x_0\\
  \Add(x_0, y+1) & = g(x_0, y, \Add(x_0, y)) =
  \Succ(\Add(x_0, y))
\end{align*}
إذن تكون $\Add$ عودية بدائية إذا كانت $f$ و$g$ كذلك. ولدينا
$f(x_0)=x_0=\Proj{1}{0}(x_0)$، ودوال الإسقاط تعد عودية بدائية، فتكون
$f$ عودية بدائية. أما $g$ فهي الدالة الثلاثية الحجج $g(x_0,y,z)$
المعرفة بـ
\[
g(x_0, y, z) = \Succ(z).
\]
ولا يخبرنا هذا بعد أن $g$ عودية بدائية، لأن $g$ و$\Succ$ ليستا الدالة
نفسها تمامًا: فـ$\Succ$ أحادية الحجة، أما $g$ فيلزم أن تكون ثلاثية
الحجج. لكن يمكننا تعريف $g$ «رسميًا» بالتركيب هكذا:
\[
g(x_0, y, z) = \Succ(\Proj{3}{2}(x_0, y, z)).
\]
ولأن $\Succ$ و$\Proj{3}{2}$ تعدان دالتين عوديتين بدائيتين، فإن $g$
كذلك، إذ يمكن تعريفها بالتركيب من دوال عودية بدائية.
\end{proof}

\begin{prop}
  \ollabel{prop:mult-pr}
  دالة الضرب $\Mult(x,y)=x\cdot y$ عودية بدائية.
\end{prop}

\begin{proof}
  تمرين.
\end{proof}

\begin{prob}
  أثبت \olref[cmp][rec][prf]{prop:mult-pr} ببيان أن تعريف $\Mult$
  بالعودية البدائية يمكن وضعه في الصورة التي تتطلبها
  \olref[cmp][rec][prf]{defn:primitive-recursion}، وببيان أن الدالتين
  المناظرتين $f$ و~$g$ عوديتان بدائيتان.
\end{prob}

\begin{ex}
إليك أول مثال عرضناه لتعريف بالعودية البدائية:
\begin{align*}
h(0) & =  1 \\
h(y+1) & =  2 \cdot h(y).
\end{align*}
لا يمكن لهذه الدالة أن تلائم الصورة المطلوبة في
\olref{defn:primitive-recursion}، لأن $k=0$. ويتضمن التعريف أيضًا
الثابتين $1$ و$2$. ولتجاوز المشكلة الأولى، ندخل حجة صورية ونعرف
الدالة~$h'$:
\begin{align*}
h'(x_0, 0) & =  f(x_0) = 1 \\
h'(x_0, y+1) & =  g(x_0, y, h'(x_0, y)) = 2 \cdot h'(x_0, y).
\end{align*}
يمكن تعريف الدالة $f(x_0)=1$ من $\Succ$ و$\Zero$ بالتركيب:
$f(x_0)=\Succ(\Zero(x_0))$. ويمكن تعريف الدالة $g$ بالتركيب من
$g'(z)=2\cdot z$ ومن دوال الإسقاط:
\begin{align*}
g(x_0, y, z) & = g'(\Proj{3}{2}(x_0, y, z))
\intertext{ويمكن بدورها تعريف $g'$ بالتركيب هكذا:}
g'(z) & = \Mult(g''(z), \Proj{1}{0}(z))
\intertext{حيث}
g''(z) & = \Succ(f(z)),
\end{align*}
و$f$ كما سبق: $f(z)=\Succ(\Zero(z))$. وبعد أن أصبحت~$h'$ لدينا،
يمكننا استعمال التركيب مرة أخرى لنجعل
$h(y)=h'(\Proj{1}{0}(y),\Proj{1}{0}(y))$. يبين هذا أن $h$ يمكن تعريفها
من الدوال الأساسية بسلسلة من التركيبات والعوديات البدائية، ومن ثم تكون
$h$ عودية بدائية.
\end{ex}

\end{document}
